\subsection{Prediction 1: Why Trade Wars Don't End When Tariffs Do}
\label{subsec:pred1-tradewar}

%%%%%%%%%%%%%%%%%%%%%%%%%%%%%%%%%%%%%%%%%%%%%%%%%%%%%%%%%%%%%%%%%%%%%%%%%%%%%%%
% CHAPTER 11.1: Narrative Treatment of Prediction 1
% Intuitive, accessible, with key data - references Appendix AU for technicals
%%%%%%%%%%%%%%%%%%%%%%%%%%%%%%%%%%%%%%%%%%%%%%%%%%%%%%%%%%%%%%%%%%%%%%%%%%%%%%%

\subsubsection{The Puzzle: Winners Who Should Be Losers}

In 2018, something strange happened in American politics.

The Trump administration imposed tariffs on \$350 billion worth of Chinese imports. 
Economists immediately calculated the costs: higher prices for consumers, disrupted 
supply chains, retaliatory tariffs on American exports. Studies confirmed the losses---
American households paid an average of \$800 more per year; farmers lost billions 
in export sales; manufacturing jobs, the supposed beneficiaries, showed no net gains 
\citep{amiti2019impact, fajgelbaum2020return}.

And yet.

The tariffs remained popular. Not just tolerated---\emph{popular}. In the very 
regions hit hardest by the economic costs, support for the trade war held steady 
or increased. Manufacturing towns that lost contracts, farming communities that 
lost markets, continued to support the policy that was, by every economic measure, 
making them poorer.

Standard economics calls this irrational. Behavioral economics calls it framing effects. 
Political science calls it tribalism.

We call it \emph{predictable}---once you understand that trade wars affect more than money.

\subsubsection{The Hidden Gains: What Economics Missed}

Consider Maria, a factory worker in Ohio. Her plant assembles components that became 
20\% more expensive after tariffs. Orders slowed. Hours were cut. Measured purely 
in dollars, she lost.

But Maria gained something too.

For twenty years, she had watched her community hollow out. Plants closed. Young 
people left. The local paper ran stories about ``forgotten Americans.'' The national 
conversation, when it mentioned places like hers at all, spoke of them as problems 
to be solved, populations to be retrained, remnants of an obsolete economy.

Then the tariffs came. And suddenly, Maria's community was at the center of 
national policy. The President spoke about ``fighting for our workers.'' News 
cameras came to town. For the first time in decades, it felt like someone was 
listening.

This isn't just psychology. It's measurable. Our framework captures it through 
the distinction between financial welfare ($U^F$) and social welfare ($U^S$).

\subsubsection{The Numbers: What We Found}

Using the $\Gamma$-matrix estimated in Chapter~10, we can quantify what Maria 
experienced. (The full technical analysis is in Appendix AU, Section~\ref{subsec:pred-01-calc}.)

The tariffs reduced market integration---what we call $\Psi_M$ in our context 
framework. This had two effects:

\begin{center}
\renewcommand{\arraystretch}{1.3}
\begin{tabular}{lcc}
\hline
\textbf{Welfare Dimension} & \textbf{Effect} & \textbf{Magnitude} \\
\hline
Financial ($U^F$) & Negative & $-0.19$ per 10pp tariff \\
Social ($U^S$) & \textbf{Positive} & $+0.06$ per 10pp tariff \\
\hline
\end{tabular}
\end{center}

The financial loss is what economists measured. The social gain is what they missed.

How can reducing market integration \emph{increase} social welfare? Because 
$\gamma_{S,M} = -0.18$---the coefficient linking market scope to social welfare 
is \emph{negative}. More global markets mean more competition, more disruption, 
more sense of being at the mercy of distant forces. Less global markets, 
paradoxically, can strengthen local community bonds.

This isn't our speculation. It's what the data show.

\subsubsection{Testing the Prediction}

If our framework is right, we should observe:
\begin{enumerate}
    \item Financial harm in tariff-exposed areas (standard prediction)
    \item Social capital \emph{gains} in those same areas (novel prediction)
    \item The gains partially offsetting the losses in overall wellbeing
\end{enumerate}

We tested this using county-level data from the 2018--2022 period.

\paragraph{What We Found on Financial Effects.}
Exactly what standard economics predicted. Counties with higher tariff exposure 
showed lower income growth, higher unemployment, more poverty. No surprises here.

\paragraph{What We Found on Social Effects.}
Here's where it gets interesting.

\begin{quote}
Counties with higher tariff exposure showed \emph{increased} community belonging, 
\emph{higher} organizational membership, and \emph{greater} trust in neighbors---
even as their financial outcomes declined.
\end{quote}

\begin{center}
\renewcommand{\arraystretch}{1.2}
\begin{tabular}{lcc}
\hline
\textbf{Social Outcome} & \textbf{Effect} & \textbf{Significant?} \\
\hline
Community belonging index & $+3.1\%$ & Yes \\
Local organization membership & $+1.8\%$ & Yes \\
Trust in neighbors & $+2.4\%$ & Yes \\
\hline
\end{tabular}
\end{center}

The tariffs made these communities poorer but more cohesive.

\paragraph{Life Satisfaction: The Net Effect.}
What about overall wellbeing? Here the effects were mixed. Life satisfaction 
showed a small, statistically insignificant decline. The social gains didn't 
fully offset the financial losses---but they significantly \emph{buffered} them.

This explains the political puzzle. People weren't ignoring the economic costs. 
They were weighing them against social and identity gains that economists 
weren't measuring.

\subsubsection{The Hysteresis Problem: Why It Doesn't End}

But here's the deeper problem. Trade wars change more than welfare levels. 
They change the \emph{context}---and context changes don't reverse.

In January 2020, the US and China signed the Phase One trade deal. Tariffs 
were partially reduced. If standard economics were right, the effects should 
have begun reversing.

They didn't.

Supply chains that had restructured during the trade war stayed restructured. 
Companies that had found alternative suppliers didn't switch back. The 
distrust between trading partners didn't heal.

We quantified this persistence. Two years after the Phase One deal:
\begin{itemize}
    \item 72\% of supply chain changes were still in place
    \item 58\% of firms maintained their new sourcing arrangements
    \item International business trust remained at trade-war levels
\end{itemize}

This is what we call \textbf{hysteresis}: the tendency of effects to persist 
after their cause is removed. It happens because trade wars don't just 
change prices---they change the underlying context in which economies operate.

\paragraph{The Formal Statement.}

We can express this precisely. Let $\Psi$ represent the economic context---
the level of market integration, institutional trust, and social cohesion. 
Trade wars shift $\Psi$:

\begin{equation}
\Psi_{after} = \Psi_{before} + \phi \cdot \text{(tariff shock)}
\end{equation}

But removing tariffs doesn't restore $\Psi$:

\begin{equation}
\lim_{t \to \infty} \Psi_t \neq \Psi_{before} \quad \text{even after tariffs removed}
\end{equation}

About 45\% of the context shift becomes permanent. The global trading system 
doesn't return to its pre-war state. It settles into a new equilibrium---
one with lower integration, lower trust, and fundamentally different 
economic relationships.

\subsubsection{What This Means for Policy}

The implications are sobering.

\paragraph{For Trade Negotiators.}
You cannot assume that removing tariffs will restore the status quo. 
Once a trade war begins, you're negotiating the terms of a new relationship, 
not returning to the old one. The question isn't ``how do we get back to 
where we were?'' but ``what new equilibrium can we reach?''

\paragraph{For Domestic Policymakers.}
Financial compensation alone won't address the effects of trade disruption. 
If communities gain social cohesion and identity validation from trade 
protection, offering them job retraining and income support misses 
the point. You're compensating for the wrong loss.

More effective approaches might include:
\begin{itemize}
    \item Recognizing and honoring community identity (not just economic adjustment)
    \item Building new sources of social cohesion that don't depend on trade barriers
    \item Gradual transitions that preserve social bonds
\end{itemize}

\paragraph{For International Institutions.}
The WTO and other trade institutions face a challenge beyond adjudicating 
disputes. Every trade war erodes \emph{institutional context}---the shared 
understanding that rules will be followed, that disputes will be resolved 
through established processes, that the system can be trusted.

This erosion is partially irreversible. Each trade conflict makes the next 
one more likely, because the institutions that prevent conflicts are 
themselves damaged by conflicts.

\subsubsection{The Deeper Pattern}

Trade wars illustrate a general principle of the framework:

\begin{quote}
\textbf{Policies that change context have effects that outlast the policies themselves.}
\end{quote}

A tariff is a price change. A trade war is a context change. Price changes 
can be reversed. Context changes---to trust, to institutions, to social 
relationships---cannot be fully undone.

This applies far beyond trade:
\begin{itemize}
    \item A financial crisis changes the context of risk perception
    \item A pandemic changes the context of social interaction
    \item A political realignment changes the context of civic engagement
\end{itemize}

In each case, the effects persist after the original cause is gone. 
The world moves to a new equilibrium, not back to the old one.

\subsubsection{Summary: Prediction 1}

\begin{tcolorbox}[colback=gray!5!white, colframe=gray!75!black, title=Prediction 1: Trade War Hysteresis]
\textbf{What standard theory predicts:}
\begin{itemize}
    \item Trade wars reduce welfare (primarily financial)
    \item Effects reverse when tariffs are removed
    \item All regions should oppose welfare-reducing policies
\end{itemize}

\textbf{What the framework predicts:}
\begin{itemize}
    \item Trade wars reduce financial welfare but may increase social welfare
    \item Effects persist after tariffs are removed (hysteresis)
    \item Regions with high social/identity weights may support trade wars
\end{itemize}

\textbf{What the data show:}
\begin{itemize}
    \item Financial effects: $-0.17$ (predicted: $-0.19$) ✓
    \item Social effects: $+0.04$ (predicted: $+0.06$) ✓
    \item Persistence: 52\% at 2 years (predicted: 45\%) ✓
\end{itemize}

\textbf{Bottom line:} Trade wars don't end when tariffs end. The context 
changes they produce---to trust, institutions, and social relationships---
are partially permanent. Policy must account for this hysteresis.
\end{tcolorbox}

\vspace{0.5em}
\noindent\textit{Technical details: See Appendix AU, Section~\ref{subsec:pred-01-calc} for the full formal model, derivation of $\beta$ coefficients, empirical validation methodology, and dynamic simulation results.}

%%%%%%%%%%%%%%%%%%%%%%%%%%%%%%%%%%%%%%%%%%%%%%%%%%%%%%%%%%%%%%%%%%%%%%%%%%%%%%%
